% AY2023 材料力学 〔II〕（転記）
% 出典: 04_材料力学/original/04_材料力学/2023/2023za.pdf
% 要約: 点Bで固定された長さ2Lの丸棒(直径d)。単位長さ当たりf0の等分布荷重が作用,
%       自由端Aにねじりトルク TA が作用。縦弾性係数E, 横弾性係数G。C は中点(Aから L)。
%       (1) 固定端B支持力RB・支持モーメントMB・支持トルクTB(二重矢印), (2) AからxのM,T,
%       (3) 点Cの上面のσx,σy,τxy, (4) 自由端AたわみδA・ねじれ角φ(たわみ基礎式),
%       (5) カスチリアーノの定理でAのたわみ導出。

\section*{〔II〕}

図2のように, 点 B で固定された長さ $2L$ の丸棒（直径 $d$）に, 単位長さ当たり $f_0$ の
等分布荷重が作用するとともに, 自由端 A にねじりトルク $T_A$ が作用している.
次の問い (1)\,$\sim$\,(5) に答えよ. 丸棒の縦弾性係数を $E$, 横弾性係数を $G$ とする.

\begin{center}
\begin{tikzpicture}[>=Latex]
  % 固定壁（右, B）
  \fill[pattern=north east lines] (8,-0.7) rectangle (8.5,0.9);
  \draw (8,-0.7) -- (8,0.9);
  % 丸棒（A=0 → B=8, 2L）
  \draw[line width=1.2pt] (0,0.1) -- (8,0.1);
  \draw[line width=1.2pt] (0,-0.1) -- (8,-0.1);
  % 等分布荷重 f0（全長, 下向き）
  \draw (0,0.95) -- (8,0.95);
  \foreach \x in {0.3,0.8,...,7.8} {\draw[->] (\x,0.95) -- (\x,0.18);}
  \node at (4.0,1.2) {$f_0$};
  % ねじりトルク TA（自由端A, 二重矢印, 右向き）
  \draw[-{Latex[length=2.2mm]}] (-1.0,0.0) -- (0.0,0.0);
  \draw[-{Latex[length=2.2mm]}] (-1.0,0.16) -- (-0.18,0.16);
  \node[left] at (-1.05,0.08) {$T_A$};
  % z, x 軸
  \draw[->] (0,0) -- (0,1.7) node[left]{$z$};
  \draw[->] (0,0) -- (8.8,0) node[right]{$x$};
  % 節点
  \node[below] at (0,-0.15) {A}; \node[below] at (4,-0.15) {C}; \node[below] at (8,-0.15) {B};
  \fill (4,0) circle (1.4pt);
  % 寸法 L, L
  \draw[<->] (0,1.5) -- (4,1.5) node[midway,fill=white]{$L$};
  \draw[<->] (4,1.5) -- (8,1.5) node[midway,fill=white]{$L$};
\end{tikzpicture}

\vspace{1mm}
図2
\end{center}

\begin{enumerate}[label=(\arabic*)]
  \item 固定端 B での支持力 $R_B$, 支持モーメント $M_B$, 支持トルク $T_B$ を図示して求めよ.
        なお, トルク表記については図2のように二重の矢印で示すこと.

  \item 自由端 A から $x$ の位置にある断面での曲げモーメント $M$ とねじりトルク $T$ を求めよ.

  \item 点 C の上面における垂直応力 $\sigma_x$, $\sigma_y$, せん断応力 $\tau_{xy}$ を求めよ.

  \item 自由端 A でのたわみ $\delta_A$ と丸棒のねじれ角 $\phi$ を求めよ.
        なお, たわみはたわみの基礎式から求めよ.

  \item カスチリアーノの定理を用いて, 自由端 A のたわみを求める導出過程を示せ.
\end{enumerate}
